% Section 3 -- Architecture. Pages target ~2.0.
% Defends C2 (manager-worker simple beats debate).
\section{Architecture}
\label{sec:arch}

The system has three persistent objects: a \emph{manager} LLM session, a
\emph{worker} LLM session per paper, and a git \emph{worktree} per worker.
Everything else is short-lived and either re-derived from these objects on
demand or written into an append-only audit log. We describe the loop
(\S\ref{sec:arch-loop}), the operational lever
(\S\ref{sec:arch-mgr}), and the contract pattern that allowed four sub-agent
modules to ship in parallel (\S\ref{sec:arch-contract}).

\subsection{Manager-Worker Adversarial Loop}
\label{sec:arch-loop}

The manager LLM is a long-lived chat session with a paranoid-reviewer
persona; the worker LLM is a long-lived chat session with a paper-writer
persona; both are persisted across process restarts via the codex CLI's
resume mechanism, \texttt{codex exec resume <session\_id>}, with the
session id captured from the first turn's stdout
(\S\ref{sec:arch-contract}). The codex-only invariant -- the in-loop
manager, worker, and reviewer-sim all run via \texttt{codex} with
different \texttt{--profile} keys -- removes the cross-CLI brittleness
that an earlier prototype carried (cursor-agent has been removed from
the loop; the binary is unavailable on remote GPU hosts and tokens are
not portable across platforms).
A scheduler (\texttt{launchd},
\texttt{cron}, or \texttt{nohup} loop) wakes the manager every five minutes
with a single token of input -- ``\texttt{tick now (<iso8601>)}'' -- so the
manager runs exactly when poked rather than busy-polling.
Algorithm~\ref{alg:tick} reproduces the manager's standard tick procedure
verbatim from its persona prompt.\footnote{Verbatim because we discovered
during dogfood that paraphrasing the procedure in the paper but leaving
the prompt unchanged was itself a hard-rule violation candidate (HR-J in
\S\ref{sec:discipline-rules}).}

\begin{algorithm}[t]
\small
\caption{Manager tick loop (one invocation per scheduler kick).}
\label{alg:tick}
\begin{algorithmic}[1]
\State \texttt{state} $\gets$ \texttt{mgr status <wid>} \Comment{worker run + worktree git status}
\If{\texttt{state.run.status} is \texttt{RUNNING} and \texttt{elapsed} $<$ \texttt{max\_run\_seconds}}
    \State \textbf{return} \texttt{\{action: skip, reason: ``in progress''\}}
\ElsIf{\texttt{state.run.status} is \texttt{RUNNING} and \texttt{elapsed} $\geq$ \texttt{max\_run\_seconds}}
    \State \texttt{out} $\gets$ \texttt{mgr output <wid>}; decide hung vs.\ legitimately long
    \State on hung: \texttt{mgr cancel}, \texttt{mgr escalate}, \textbf{return}
\EndIf
\If{\texttt{state.run.id} $=$ \texttt{state.last\_processed\_run\_id}}
    \State \textbf{goto} step~\ref{algline:newrun} \Comment{no new evidence, dispatch next task}
\EndIf
\State \texttt{commits} $\gets$ \texttt{mgr commits <wid> -n 10}
\If{any commit subject starts with \texttt{[BLOCKED-DECISION]}}
    \State \texttt{mgr escalate <wid>}; \textbf{return} \texttt{\{action: escalate\}}
\EndIf
\State \texttt{diff} $\gets$ \texttt{mgr diff <wid>}; \texttt{rules} $\gets$ \texttt{mgr rules <wid>}
\State Check each hard rule against \texttt{diff} citing concrete file:line evidence
\If{paper-claim file (abstract / intro / conclusion) touched \textbf{or} diff $\geq 2$ files}
    \State \label{alg:review-sim-guard} \texttt{mgr review-sim <wid>} \Comment{cross-model NeurIPS-reviewer judgement}
\EndIf
\If{any rule violated and not auto-correctable}
    \State \texttt{mgr escalate <wid>}; \textbf{return} \texttt{\{action: escalate\}}
\ElsIf{rule violated and auto-correctable}
    \State \texttt{mgr cancel}; \texttt{mgr start <wid> -p ``<corrective prompt>''}
\Else
    \State \label{algline:newrun} \texttt{mgr start <wid> -p ``<next-step prompt>''}
\EndIf
\State \texttt{mgr mark-processed <wid> --run-id <id> --verdict <v> --summary <s>}
\State Emit single-line JSON audit \texttt{\{tick\_at, worker, action, summary, run\_id\}}
\end{algorithmic}
\end{algorithm}

Two structural properties distinguish this loop from existing multi-agent
paper-writing systems. First, manager and worker have \emph{adversarial
roles by construction}: the manager's persona prompt states that it is
not present to help the worker but to prevent wrong-but-plausible
commits, and all communication between them goes through the manager's
reading of worker-produced git commits rather than direct dialog.
Second, the manager is \emph{stateless within a tick} but persistent
\emph{across} ticks: each wake-up loads only the small \texttt{state}
JSON and the manager's own chat history, never the worker's, so the
manager cannot be talked out of a verdict by a verbose worker.

\subsection{The mgr CLI as the only lever}
\label{sec:arch-mgr}

The manager LLM is given exactly one tool, a Python CLI named \texttt{mgr}
with thirteen atomic operations: four \emph{read} operations
(\texttt{workers}, \texttt{status}, \texttt{commits}, \texttt{diff}) for
inspecting worker state and worktree git diffs; three \emph{introspection}
operations (\texttt{output}, \texttt{rules}, \texttt{audit},
\texttt{log}, \texttt{escalations}) for reading worker stdout, the
worker's own hard-rules file, the manager's tick audit history, the
event stream, and the pending escalation queue; three \emph{control}
operations (\texttt{start}, \texttt{cancel}, \texttt{escalate}) for
forking / killing worker processes and marking items for human
decision; and one \emph{state} operation
(\texttt{mark-processed}) for recording a verdict against a run id.
The four v2 sub-agent commands of \S\ref{sec:discipline}
(\texttt{review-sim}, \texttt{lit-review}, \texttt{sentinel-tick},
\texttt{backend-test}) are wired in identically. The manager cannot edit
files directly, run \texttt{git push}, change worker prompts, or call
any LLM other than itself; every action it wants to take must factor
through one of these atoms. The Python implementation of the CLI is
\(789\) lines (out of \(5{,}228\) total LOC for the released tree in Python
plus Bash plus the sub-agent contract Markdown, counted on 2026-05-04 by
\texttt{wc -l mgr lib/*.py sentinel.py scripts/*.py scripts/*.sh *.sh docs/sub\_agents\_contract.md} at the public-release commit cited in
Appendix~\ref{app:repo}); this
restriction is deliberate -- the manager's job is operational reasoning,
not authoring, and constraining its action space prevents the most
common failure mode of prior systems where a critic agent quietly
rewrites the artifact it is supposed to be critiquing.

The worker, by contrast, has free shell access inside its git worktree --
it can write any file, run \texttt{pdflatex}, query Semantic Scholar, and
push to its own isolated branch -- but it cannot read or write any state
outside its worktree, including the manager's persona, other workers'
worktrees, or the audit logs. This asymmetry is what makes the loop
adversarial without being deadlocked: the worker has the authoring
privilege, the manager has the verdict privilege, and each can do exactly
the things the other cannot.

Worker isolation uses standard \texttt{git worktree}: multiple
workers share one repository clone, commit to disjoint branches, do
not see each other's edits, and never push to remote (commits
accumulate on the local worker branch until a human merges them via
a non-fast-forward merge). All audit logs are append-only JSONL
written by \texttt{mgr} and form the immutable ground truth that
\S\ref{sec:case} mines for empirical claims where those artefacts are
available (git history and verifier transcripts; default developer
checkouts \texttt{gitignore} \texttt{state/}, so runtime JSONL is not
claimed to ship in full in the public OSS tree -- see~\S\ref{sec:case-setup}).

\subsection{Sub-agent Contract Design}
\label{sec:arch-contract}

The four v2 sub-agents (\texttt{reviewer\_sim}, \texttt{lit\_review},
\texttt{sentinel}, \texttt{codex\_worker}; \S\ref{sec:discipline}) were
implemented in parallel by four LLM-driven implementation streams in
under thirty wall-clock minutes -- a wall-clock improvement of
approximately two orders of magnitude over our initial three-to-four
work-day estimate. The mechanism is what we call a \emph{sub-agent
contract}: a single Markdown file
(\texttt{docs/sub\_agents\_contract.md}, 593 lines) that fully specifies,
for each module to be shipped in parallel,

(i) the exact set of files the implementer is allowed to create or
modify, plus an explicit ``do not touch'' list naming all sibling
sub-agents' files; (ii) the public Python API (function signatures,
return-value dataclasses, exceptions raised); (iii) the JSON / Markdown
output schema written to
\texttt{state/<module>/<wid>/<run\_id>.\{json,md\}}; (iv) a self-test
(\texttt{python -m lib.<module> --self-test}) that must run
hermetically in under ten seconds with zero network calls; and (v) the
parent \texttt{mgr} subcommand wiring (argparse arguments, return
codes, stdout format).

Because every public surface that crosses module boundaries is named in
the contract, the implementation streams can run completely in parallel
with zero shared mutable state and zero coordination overhead. We
verified the resulting modules contain zero cross-imports between
sibling sub-agents and integrate cleanly into the parent \texttt{mgr}
CLI through a single hand-authored 80-line factory function. The
thirteen-pass smoke test (\texttt{scripts/smoke\_test\_sub\_agents.sh})
runs all four self-tests plus the parent argparse wiring in under two
seconds and was the gating criterion for the integration commit.

We position contract-first parallel agent shipping as a methodological
contribution in its own right (independent of the discipline mechanisms
it enables) and give the full contract template in
Appendix~\ref{app:contract-template}.
